\documentclass[11pt]{article}

\usepackage[margin=1.35cm]{geometry}
\usepackage{amsmath,amssymb}
\usepackage{enumitem}
\usepackage{array}

\setlength{\columnsep}{0.65cm}
\setlength{\parindent}{0pt}
\setlist[enumerate]{leftmargin=*,itemsep=0.35em,topsep=0.25em}

% Multiple-choice bubble
\newcommand{\choice}[2]{%
    \par\noindent
    \(\bigcirc\)\ \textbf{#1.}\ #2
}

% Compact answer-sheet bubble row
% Reduced font size so A--D fit on one line in each table cell.
\newcommand{\bubblecell}[1]{%
    {\scriptsize
    \textbf{#1}\;
    \(\bigcirc\)A\;
    \(\bigcirc\)B\;
    \(\bigcirc\)C\;
    \(\bigcirc\)D}%
}

\begin{document}

\begin{center}
{\Large\bfseries COSC 120: Introduction to Python Programming}\\[2pt]
{\large\bfseries In-Class Midterm 2 Examination -- Version 3}\\[3pt]
Prof. Jordan C. Hanson
\end{center}

\vspace{0.1cm}

\noindent
\textbf{Name:} \hrulefill
\hspace{0.5cm}
\textbf{Student ID:}
\fbox{\rule{0pt}{1.8ex}\hspace{4cm}}

\vspace{0.2cm}

\noindent
\textbf{Instructions:}
Choose the best answer for each question.
Fill in the circle next to exactly one answer, A--D.
Each question is worth 4 points.

\vspace{0.15cm}
\hrule
\vspace{0.2cm}

\begin{center}
{\large\bfseries Answer Sheet}
\end{center}

\small
\renewcommand{\arraystretch}{1.3}

\begin{center}
\begin{tabular}{
|m{0.18\textwidth}
|m{0.18\textwidth}
|m{0.18\textwidth}
|m{0.18\textwidth}
|m{0.18\textwidth}|}
\hline
\bubblecell{1} &
\bubblecell{2} &
\bubblecell{3} &
\bubblecell{4} &
\bubblecell{5} \\ \hline

\bubblecell{6} &
\bubblecell{7} &
\bubblecell{8} &
\bubblecell{9} &
\bubblecell{10} \\ \hline

\bubblecell{11} &
\bubblecell{12} &
\bubblecell{13} &
\bubblecell{14} &
\bubblecell{15} \\ \hline

\bubblecell{16} &
\bubblecell{17} &
\bubblecell{18} &
\bubblecell{19} &
\bubblecell{20} \\ \hline

\bubblecell{21} &
\bubblecell{22} &
\bubblecell{23} &
\bubblecell{24} &
\bubblecell{25} \\ \hline
\end{tabular}
\end{center}

\normalsize

\vspace{0.15cm}
\hrule
\vspace{0.25cm}

\twocolumn

%=========================================================
\section*{Dictionaries and Comprehensions}

\begin{enumerate}

%---------------------------------------------------------
% Question 1
%---------------------------------------------------------
\item Consider the dictionary

\begin{verbatim}
students = {
    2601001: "Ana",
    2601002: "Ben",
    2601003: "Chloe"
}
\end{verbatim}

What is produced by

\begin{verbatim}
students[2601002]
\end{verbatim}

\choice{A}{Ana}
\choice{B}{Ben}
\choice{C}{2601002}
\choice{D}{An error occurs}

%---------------------------------------------------------
% Question 2
%---------------------------------------------------------
\item Suppose the key \texttt{2601999} is not in \texttt{students}. What is returned by

\begin{verbatim}
students.get(2601999)
\end{verbatim}

\choice{A}{0}
\choice{B}{False}
\choice{C}{None}
\choice{D}{The program always stops with an error}

%---------------------------------------------------------
% Question 3
%---------------------------------------------------------
\item Consider

\begin{verbatim}
names = ["Ana", "Ben", "Chloe",
         "Diego", "Elena"]
first_id = 2601001
students = {
    first_id + i: name
    for i, name in enumerate(names)
}
\end{verbatim}

Which key is associated with \texttt{"Diego"}?

\choice{A}{2601001}
\choice{B}{2601002}
\choice{C}{2601003}
\choice{D}{2601004}

%---------------------------------------------------------
% Question 4
%---------------------------------------------------------
\item Which statement best describes the difference between these two expressions?

\begin{verbatim}
[first_id + i
 for i in range(len(names))]

{first_id + i: name
 for i, name in enumerate(names)}
\end{verbatim}

\choice{A}{The first creates a list of IDs; the second creates ID--name pairs in a dictionary}
\choice{B}{Both create dictionaries}
\choice{C}{The first creates a dictionary; the second creates a list}
\choice{D}{Both create exactly the same object}

%---------------------------------------------------------
% Question 5
%---------------------------------------------------------
\item In our lookup-timing exercise, why did we use

\begin{verbatim}
time.perf_counter()
\end{verbatim}

\choice{A}{To sort the student records}
\choice{B}{To measure short execution times with a high-resolution timer}
\choice{C}{To count how many dictionary keys exist}
\choice{D}{To pause execution for a fixed number of seconds}

%---------------------------------------------------------
% Question 6
%---------------------------------------------------------
\item A student ID is near the end of a nested list containing \(N\) records. Which statement best describes our timing experiment as \(N\) becomes large?

\choice{A}{Both nested-list and dictionary lookup times must grow in exactly the same way}
\choice{B}{Dictionary lookup becomes impossible once \(N>1000\)}
\choice{C}{Nested-list lookup generally grows with \(N\), while dictionary lookup changes much less}
\choice{D}{Nested-list lookup becomes constant-time because the list is nested}

\end{enumerate}

%=========================================================
\section*{Classes and Objects}

\begin{enumerate}
\setcounter{enumi}{6}

%---------------------------------------------------------
% Question 7
%---------------------------------------------------------
\item In a class definition, what is the main purpose of

\begin{verbatim}
def __init__(self, ...):
\end{verbatim}

\choice{A}{To initialize a new object's attributes}
\choice{B}{To delete an object}
\choice{C}{To import another module}
\choice{D}{To begin a recursive call}

%---------------------------------------------------------
% Question 8
%---------------------------------------------------------
\item Inside an instance method, what does \texttt{self} refer to?

\choice{A}{The class file on disk}
\choice{B}{Every object that has ever been created}
\choice{C}{The Python interpreter}
\choice{D}{The particular object on which the method was called}

%---------------------------------------------------------
% Question 9
%---------------------------------------------------------
\item Consider the complete class fragment

\begin{verbatim}
class Particle2D:
    def __init__(self, x, y, vx, vy):
        self.x = x
        self.y = y
        self.vx = vx
        self.vy = vy

    def update(self, dt):
        self.x += self.vx * dt
        self.y += self.vy * dt

p = Particle2D(1, 2, 3, -1)
p.update(0.5)
print(p.x, p.vx)
\end{verbatim}

What is printed?

\choice{A}{\texttt{2.5 3}}
\choice{B}{\texttt{2.5 1.5}}
\choice{C}{\texttt{1 3}}
\choice{D}{\texttt{4 -1}}

% Question 10
%---------------------------------------------------------
\item Using the same \texttt{Particle2D} class, consider

\begin{verbatim}
p1 = Particle2D(0, 0, 2, 1)
p2 = Particle2D(0, 0, -1, 3)

p1.update(1.0)
print(p2.x, p2.y)
\end{verbatim}

What is printed, and why?

\choice{A}{\texttt{2 1}, because all instances share the same position}
\choice{B}{\texttt{-1 3}, because creating \texttt{p2} automatically updates it}
\choice{C}{\texttt{0 0}, because \texttt{update()} was called only on \texttt{p1}}
\choice{D}{An error occurs because two instances cannot exist at once}

% Question 11
%---------------------------------------------------------
\item A student writes the following method inside \\ \texttt{Particle2D}:

\begin{verbatim}
def speed(self):
    return math.sqrt(vx**2 + vy**2)
\end{verbatim}

Why will this method fail when called as \texttt{p.speed()}?

\choice{A}{The method should use \texttt{self.vx} and \texttt{self.vy}}
\choice{B}{\texttt{math.sqrt()} cannot be used inside a class}
\choice{C}{Instance methods cannot return values}
\choice{D}{\texttt{speed()} must change \texttt{x} and \texttt{y}}

% Question 12
%---------------------------------------------------------
\item Let

\begin{verbatim}
z1 = ComplexNumber(3, 2)
z2 = ComplexNumber(5, -4)
z3 = z1 + z2
\end{verbatim}

If \texttt{\_\_add\_\_()} correctly implements complex-number addition, and \texttt{z3.real} and \texttt{z3.imag} produce the real and imaginary parts of \texttt{z3}, what are \texttt{z3.real} and \texttt{z3.imag}?

\choice{A}{\texttt{15, -8}}
\choice{B}{\texttt{-2, 6}}
\choice{C}{\texttt{8, 6}}
\choice{D}{\texttt{8, -2}}

%---------------------------------------------------------
% Question 13
%---------------------------------------------------------
\item Consider the method

\begin{verbatim}
def __mul__(self, other):
    real = (self.real * other.real
            - self.imag * other.imag)
    imag = (self.real * other.imag
            + self.imag * other.real)
    return ComplexNumber(real, imag)
\end{verbatim}

Which statement best describes what happens when Python evaluates
\texttt{z3 = z1 * z2}?

\choice{A}{A new \texttt{ComplexNumber} object is created and assigned to \texttt{z3}}
\choice{B}{\texttt{z1} is permanently changed into the product}
\choice{C}{The method returns a two-element list}
\choice{D}{Python ignores \texttt{\_\_mul\_\_()} because \texttt{z1} and \texttt{z2} are objects}

\end{enumerate}

%=========================================================
\section*{Recursion and Binary Search}

\begin{enumerate}
\setcounter{enumi}{13}

%---------------------------------------------------------
% Question 14
%---------------------------------------------------------
\item Which statement best defines a \textbf{base case} in a recursive function?

\choice{A}{The first line of the function}
\choice{B}{The case that makes the largest recursive call}
\choice{C}{A stopping case that returns without making another recursive call}
\choice{D}{A loop that always runs at least once}

%---------------------------------------------------------
% Question 15
%---------------------------------------------------------
\item Consider the recursive pattern from class:

\begin{verbatim}
def all_positive(values):
    if len(values) == 0:
        return True
    x = values.pop()
    if x <= 0:
        return False
    return all_positive(values)
\end{verbatim}

What is returned for

\begin{verbatim}
all_positive([6, 2, -9, 4])
\end{verbatim}

\choice{A}{False}
\choice{B}{True}
\choice{C}{-9}
\choice{D}{None}

%---------------------------------------------------------
% Question 16
%---------------------------------------------------------
\item Consider

\begin{verbatim}
def fact(n):
    if n <= 1:
        return 1
    return n * fact(n - 1)
\end{verbatim}

When evaluating \texttt{fact(4)}, which sequence of function calls is made before the base case returns?

\choice{A}{\texttt{fact(4), fact(3), fact(2), fact(1)}}
\choice{B}{\texttt{fact(4), fact(5), fact(6), ...}}
\choice{C}{\texttt{fact(4)} only}
\choice{D}{\texttt{fact(1), fact(2), fact(3), fact(4)}}

% Question 17
%---------------------------------------------------------
\item A student writes

\begin{verbatim}
def fact(n):
    if n <= 1:
        return 1
    return n * fact(n)
\end{verbatim}

Why does this implementation fail for \texttt{fact(5)}?

\choice{A}{The base case should return 0}
\choice{B}{The recursive call uses the same \texttt{n}, so it never moves toward the base case}
\choice{C}{Multiplication is not allowed in a recursive function}
\choice{D}{Recursive functions must contain a \texttt{for} loop}

% Question 18
%---------------------------------------------------------
\item Consider this part of recursive binary search:

\begin{verbatim}
mid = (low + high) // 2

if values[mid] == target:
    return mid
if target < values[mid]:
    return binary_search(
        values, target, low, mid - 1)
return binary_search(
    values, target, mid + 1, high)
\end{verbatim}

For

\begin{verbatim}
values = [3, 7, 11, 18, 23, 29, 34]
target = 29
low = 0
high = 6
\end{verbatim}

the first midpoint is index 3, whose value is 18. What recursive call is made next?

\choice{A}{\texttt{binary\_search(values, 29, 0, 2)}}
\choice{B}{\texttt{binary\_search(values, 29, 4, 6)}}
\choice{C}{\texttt{binary\_search(values, 29, 0, 6)}}
\choice{D}{The function immediately returns \texttt{-1}}

% Question 19
%---------------------------------------------------------
\item In our recursive binary-search function, what does

\begin{verbatim}
if low > high:
    return -1
\end{verbatim}

mean?

\choice{A}{There are no indices left to search, so the target is absent}
\choice{B}{The target must be at index 0}
\choice{C}{The list should be reversed}
\choice{D}{The target has been found at \texttt{mid}}

%---------------------------------------------------------
% Question 20
%---------------------------------------------------------
\item Why is binary search efficient on a large sorted list?

\choice{A}{It always examines every element exactly once}
\choice{B}{It converts the list into a dictionary first}
\choice{C}{It works only because recursive functions are always constant-time}
\choice{D}{Each decision discards about half of the remaining search region}

\end{enumerate}

%=========================================================
\section*{Inheritance and Vector2D}

\begin{enumerate}
\setcounter{enumi}{20}

%---------------------------------------------------------
% Question 21
%---------------------------------------------------------
\item Consider the class method

\begin{verbatim}
def __add__(self, other):
    return Vector2D(
        self.x + other.x,
        self.y + other.y)
\end{verbatim}

and the code

\begin{verbatim}
v1 = Vector2D(2, 5)
v2 = Vector2D(-1, 3)
v3 = v1 + v2
\end{verbatim}

Which statement is correct?

\choice{A}{\texttt{v3} is a \texttt{Vector2D} with components \texttt{(1, 8)}}
\choice{B}{\texttt{v3} is the integer \texttt{9}}
\choice{C}{\texttt{v1} is changed to \texttt{(-1, 3)}}
\choice{D}{Python cannot use \texttt{+} with objects}

% Question 22
%---------------------------------------------------------
\item Suppose \texttt{ComplexNumber} inherits from \texttt{Vector2D} and overrides only \texttt{\_\_add\_\_()}:

\begin{verbatim}
class ComplexNumber(Vector2D):
    def __add__(self, other):
        return ComplexNumber(
            self.x + other.x,
            self.y + other.y)
\end{verbatim}

If \texttt{Vector2D} already defines \texttt{magnitude()}, which methods are used by the following expressions?

\begin{verbatim}
z1.magnitude()
z1 + z2
\end{verbatim}

\choice{A}{Both use methods defined only in \texttt{Vector2D}}
\choice{B}{\texttt{magnitude()} is inherited from \texttt{Vector2D}; addition uses the overridden method in \texttt{ComplexNumber}}
\choice{C}{Both require methods to be rewritten inside \texttt{ComplexNumber}}
\choice{D}{\texttt{magnitude()} uses \texttt{ComplexNumber}; addition uses \texttt{Vector2D}}

% Question 23
%---------------------------------------------------------
\item Consider

\begin{verbatim}
class ComplexNumber(Vector2D):
    ...
\end{verbatim}

Which statement is correct?

\choice{A}{\texttt{ComplexNumber} is a child class that inherits from \texttt{Vector2D}}
\choice{B}{\texttt{Vector2D} inherits from \texttt{ComplexNumber}}
\choice{C}{The two classes are unrelated}
\choice{D}{Python combines the two classes into a dictionary}

%---------------------------------------------------------
% Question 24
%---------------------------------------------------------
\item In

\begin{verbatim}
def __init__(self, real, imag):
    super().__init__(real, imag)
\end{verbatim}

what does \texttt{super().\_\_init\_\_()} do?

\choice{A}{It creates an unrelated \texttt{Vector2D} object and discards it}
\choice{B}{It calls the recursive base case}
\choice{C}{It overrides every method in the parent class}
\choice{D}{It reuses the parent class constructor to initialize the inherited components}

%---------------------------------------------------------
% Question 25
%---------------------------------------------------------
\item Suppose the parent and child classes contain

\begin{verbatim}
class Vector2D:
    def __add__(self, other):
        return Vector2D(
            self.x + other.x,
            self.y + other.y)

class ComplexNumber(Vector2D):
    def __add__(self, other):
        return ComplexNumber(
            self.x + other.x,
            self.y + other.y)
\end{verbatim}

For two \texttt{ComplexNumber} objects \texttt{z1} and \texttt{z2}, what is the difference between

\begin{verbatim}
a = Vector2D.__add__(z1, z2)
b = z1 + z2
\end{verbatim}

\choice{A}{Both \texttt{a} and \texttt{b} are \texttt{Vector2D} objects}
\choice{B}{\texttt{a} is a \texttt{Vector2D}; \texttt{b} is a \texttt{ComplexNumber}}
\choice{C}{\texttt{a} is a \texttt{ComplexNumber}; \texttt{b} is a \texttt{Vector2D}}
\choice{D}{Both statements produce an error}

\end{enumerate}

\vfill

\begin{center}
\textbf{End of Exam}

Check that exactly one answer is clearly marked for each question.
\end{center}

\end{document}
